% =========================================================
%  EXAM CONFIGURATION -- the only file to edit per sitting
% =========================================================
%
% groups.tex in this folder is one line that \input's ../common/exam-layout.tex.
% That layout loads every file in ../questions/ and then reads this file.
%
% Step 1: set the exam info below
% Step 2: edit \ExamParts -- one \ExamPart per part, one \DrawGroup or
%         \DrawAll per chapter inside it
% Step 3: edit \AnswerSheetLayout -- one \SheetGrid per part
% Step 4: build (see guide/05-building.md)
%
% =========================================================

% --- Exam Info ---
\newcommand{\ExamTitle}{Final Examination 1/2026}
\newcommand{\CourseCode}{999999}
\newcommand{\CourseName}{Sample Course}
\newcommand{\CourseSection}{701--702}
\newcommand{\ExamDate}{1 December 2026}
\newcommand{\ExamDateNum}{01/12/2026}   % short form for the answer-sheet header
\newcommand{\ExamTime}{09.00 -- 11.00}
\newcommand{\NumCopies}{4}      % Distinct shuffled papers. 2-4 while testing,
                                % students + a few spares for the real print.

% 0 = every copy is as long as it needs; N = pad the question pages to page N
% so all copies have the same length.
\newcommand{\ExamPadToPage}{0}

% =========================================================
%  THE PARTS
% =========================================================
% \DrawGroup{group}{N}  N random questions, a different draw on every copy.
% \DrawAll{group}       the whole group, shuffled per copy.
% \ExamPart[cols]{title}{...}: cols = 1 or 2 (default 2) question columns.
% N must never exceed the group's size (tools/check_questions.py checks it).
%
% Bank sizes:   ch01_data 8   ch02_stats 6
% =========================================================
\newcommand{\ExamParts}{%
  \ExamPart{Part 1 --- Chapter 1: Data Basics}{%
    \DrawGroup{ch01_data}{6}%
  }%
  \ExamPart{Part 2 --- Chapter 2: Statistics}{%
    \DrawAll{ch02_stats}%
  }%
}

% =========================================================
%  ANSWER SHEET LAYOUT
% =========================================================
% \SheetGrid{columns}{which parts}. \InPart{p1} is the first \ExamPart above,
% \InPart{p2} the second, and so on. Four-choice questions fit 3 columns; questions
% with more than 5 choices need a 2-column grid of their own.
\newcommand{\AnswerSheetLayout}{%
  \SheetGrid{3}{\InPart{p1} || \InPart{p2}}%
}
